% ---------------------------------------------------------------------------
% Payoff scaling shapes cooperation in LLM agents across languages
% ---------------------------------------------------------------------------
\documentclass[11pt,a4paper]{article}

\usepackage[T1]{fontenc}
\usepackage[utf8]{inputenc}
\usepackage{lmodern}
\usepackage[margin=1in]{geometry}
\usepackage{graphicx}
\usepackage{amsmath,amssymb}
\usepackage{booktabs}
\usepackage{tabularx}
\usepackage{caption}
\usepackage{authblk}
\usepackage[numbers,sort&compress]{natbib}
\usepackage{microtype}
\usepackage{float}
\usepackage{placeins}
\usepackage{enumerate}
\usepackage{xcolor}
\usepackage[colorlinks=true,allcolors=blue!55!black]{hyperref}

\captionsetup{font=small,labelfont=bf,width=\textwidth}
\setlength{\parskip}{0pt}
\renewcommand{\baselinestretch}{1.15}
\raggedbottom

% Keep substantial figures at the top of a text page instead of creating
% sparsely populated float pages in the neutral article layout.
\setcounter{topnumber}{2}
\setcounter{bottomnumber}{2}
\renewcommand{\topfraction}{0.95}
\renewcommand{\bottomfraction}{0.80}
\renewcommand{\textfraction}{0.05}
\renewcommand{\floatpagefraction}{0.80}

\newcommand{\lam}{\lambda}
\newcommand{\ci}[2]{(95\% CI #1, #2)}

\newif\ifhighlightrevisions
\highlightrevisionstrue
\ifhighlightrevisions
  \newcommand{\rev}[1]{{\color{red}#1}}
\else
  \newcommand{\rev}[1]{#1}
\fi

% In the preamble
\usepackage{authblk}
\usepackage{hyperref}

\title{Payoff scaling shapes cooperation in LLM agents across languages}

\author[]{%
The Anh Han\textsuperscript{1\textdagger*},
Trung-Kiet Huynh\textsuperscript{2,4\textdagger},
Dao-Sy Duy-Minh\textsuperscript{2,4\textdagger},
Thanh-Bang Cao\textsuperscript{3,4},
Phong-Hao Le\textsuperscript{3,4},
Hong-Dan Nguyen\textsuperscript{3,4},
Phu-Quy Nguyen-Lam\textsuperscript{2,4},
Minh-Luan Nguyen-Vo\textsuperscript{3,4},
Hong-Phat Pham\textsuperscript{3,4},
Phu-Hoa Pham\textsuperscript{2,4},
Thien-Kim Than\textsuperscript{3,4},
Chi-Nguyen Tran\textsuperscript{3,4},
Huy Tran\textsuperscript{3,4},
Gia-Thoai Tran-Le\textsuperscript{3,4},
Alessio Buscemi\textsuperscript{5},
and Le Hong Trang\textsuperscript{3,4*}
}

\affil[]{%
\textsuperscript{1}School of Computing, Engineering and Digital Technologies,
Teesside University, Middlesbrough, United Kingdom\\
\textsuperscript{2}Faculty of Information Technology,
University of Science (HCMUS), Ho Chi Minh City, Vietnam\\
\textsuperscript{3}Faculty of Computer Science and Engineering,
Ho Chi Minh City University of Technology (HCMUT),
Ho Chi Minh City, Vietnam\\
\textsuperscript{4}Vietnam National University,
Ho Chi Minh City (VNU-HCM), Ho Chi Minh City, Vietnam\\
\textsuperscript{5}Luxembourg Institute of Science and Technology (LIST),
Luxembourg
}

\renewcommand\Authfont{\normalsize}
\renewcommand\Affilfont{\small\normalfont}
\setlength{\affilsep}{0.5em}

\date{%
\vspace{-1.0em}
\footnotesize
\textsuperscript{\textdagger}\,These authors contributed equally to this work.\\
\textsuperscript{*}\,Corresponding authors:
\href{mailto:t.han@tees.ac.uk}{t.han@tees.ac.uk}
and
\href{mailto:lhtrang@hcmut.edu.vn}{lhtrang@hcmut.edu.vn}.
}

\begin{document}

\maketitle

\begin{abstract}
Large language models (LLMs) are increasingly deployed as autonomous agents that negotiate, coordinate, and make decisions on behalf of human users. Understanding how they cooperate is therefore important for AI governance and machine behaviour research. We examine two factors: payoff magnitude and interaction language. Classical game theory predicts that multiplying all payoffs by a positive constant should not change behaviour because the underlying strategic incentives remain unchanged. We test this prediction using five widely adopted commercial LLMs playing an iterated prisoner’s dilemma across ten payoff scales and five languages \rev{(English, French, Vietnamese, Chinese and Arabic)}. Contrary to theoretical expectations, changing only the numerical scale of payoffs substantially alters cooperation. Models respond differently, often in opposite directions, meaning that averages across models can conceal important behavioural heterogeneity. Sensitivity to payoff scale also depends strongly on prompt language. Moreover, payoff scale modifies the effects of persona instructions: at sub-unit stakes\rev{, where no payoff shown exceeds 1,} all five models cooperate more under a selfish than a cooperative persona, whereas Claude and Gemini reverse this pattern at larger scales. These findings show that cooperation measures depend critically on payoff scale and language and should not be treated as fixed model traits.
%Large language models (LLMs) increasingly act as autonomous agents negotiating and deciding for human users, making their cooperative dispositions a central question for AI governance. We examine how two fundamental levers shape cooperation: the nominal magnitude of the stakes and the language in which the interaction is framed. In classical game theory, multiplying all payoffs by a positive constant is an invariant transformation that leaves preference rankings, best replies, and equilibria untouched. We show, across five widely adopted commercial LLMs playing an iterated prisoner's dilemma across ten payoff scales and five languages, this invariance is systematically violated. Rescaling stakes alters cooperation rates dramatically, driving divergent shifts across models that cancel out when averaged. Scale sensitivity also varies across languages. Payoff scale also changes persona compliance: at sub-unit stakes, all five models cooperate more under the selfish than the cooperative persona; in Claude and Gemini this effect is reversed above the sub-unit regime. Across the evaluated models, own-action conditioning is stronger than partner-action conditioning. These findings show that cooperation metrics cannot be treated as intrinsic model traits without specifying payoff scale and language.
\end{abstract}

\noindent\textbf{Keywords:} machine behaviour; large language models; iterated prisoner's dilemma; payoff scaling; utility invariance; cross-lingual behaviour; (evolutionary) game theory


\section{Introduction}\label{sec:intro}

Large language models \rev{(LLMs)} are increasingly deployed as agents that negotiate, coordinate and decide on behalf of their users, in recommendation systems, negotiation tools and multi-agent assistants~\cite{park2023generative,tessler2024ai,hammond2025multi,cook2024llm}. In these settings, behaviour emerges from strategic interaction rather than from isolated model outputs, and the agents face cooperation dilemmas with solutions  directly relevant to  safety, coordination and AI governance~\cite{dafoe2021cooperative,lu2024llms,han2026social,idowu2026mapping}. \rev{More broadly, LLM-driven generative agents have been proposed as a new route for studying complex systems and evolutionary games~\cite{lu2024llms,han2025llms,chen2025harnessing}.} A fast-growing literature measures how they behave when they do, using methods from behavioural economics and cognitive psychology~\cite{binz2023using,horton2023large,aher2023using,rilla2026recognising} and  the classical games~\cite{mei2024turing}. Equally consequential is how far that behaviour adapts to the costs and benefits at issue, since a deployed agent meets the same dilemma at many different magnitudes~\cite{hammond2025multi}.

A growing body of work evaluates language models through game-theoretic lenses, using evolutionary and repeated games, reporting systematic departures from equilibrium outcomes, persistent cooperative biases and sensitivity to contextual factors such as language and incentives~\cite{buscemi2025fairgame,sun2025gtllm_survey,mao2025alympics,pal2026strategiescooperationdefectionlarge,willis2025will,fan2024can,pires2025large,fu2026optimal,xie2026evolutionary}. The prisoner's dilemma has been a well adopted framework ~\cite{axelrod1981evolution,axelrod1984evolution,nowak1992tit,nowak1993strategy,nowak2006five,pressDyson2012}. Our study is closest to FAIRGAME, which introduced a game-theoretic evaluation framework and reported payoff and language sensitivity in repeated games~\cite{buscemi2025fairgame}. We retain its no-communication prisoner's dilemma protocol so that the comparison is anchored to the same task, while adding a positive-rescaling null spanning ten scales, five languages and five model endpoints. This makes our contribution a test of invariance and heterogeneity, rather than an evaluation of LLM behaviours in a new strategic interaction.

Prior work showed that the LLM models differ in how readily they retaliate~\cite{akata2025} or forgive~\cite{fontana2025nicer}, how much weight they give the game rather than its framing  ~\cite{loreheydari2024strategic} and how far an instructed disposition carries into conditioned reciprocity~\cite{phelps2023investigating}. These studies, however, share two largely unremarked reporting conventions. First, they characterise a model using a cooperation rate, measured under a single payoff matrix. Second, they report what agents did, rather than the decision rule that generated the observed behaviour. This makes it difficult to tell whether a reported cooperation rate is a stable property of the model, an artefact of the particular payoff scale and framing, or the aggregate trace of a different strategic rule.

Evaluating strategic behaviour at the level required for alignment and governance therefore requires more than aggregate outcomes. Drawing on behavioural and evolutionary game theory~\cite{sigmund:2010bo}, a behavioural intention can be conceptualised as an agent's \emph{strategy}: a decision rule mapping an interaction history to the next action~\cite{han2013intention,diStefanoIntention2023,Han2012ALIFEjournal,fujimoto2019functional}. The canonical strategies of repeated games, e.g. Always Cooperate, Always Defect, Tit-for-Tat and Win-Stay-Lose-Shift~\cite{Axelrod1980Effective,sigmund:2010bo}, provide an interpretable vocabulary for such intentions, and prior work has shown that they can be inferred from noisy behavioural trajectories by supervised learning~\cite{han2013intention,diStefanoIntention2023} and extended by probabilistic models to mixed or stochastic play~\cite{inferenceStrategies}. The inference is not free. Separating the memory-one rules from observed play has been shown to need a long run of observations even where the resident rule is fixed by construction~\cite{kim2021win,hilbe2013memory}, which is why the read-out used here is reported as a resemblance, with its provenance beside it (Section \ref{sec:readout}).

In parallel, a language model's behaviour varies with the language in which the prompt is written ~\cite{buscemi2025fairgame}. Strategic reasoning, risk sensitivity and cooperation patterns depend on linguistic framing, sometimes by margins comparable to the differences between models~\cite{loreheydari2024strategic,buscemi2025llms,buscemi2025strategic}. Yet it remains unclear whether the strategies that agents appear to play are themselves language-dependent, and whether such dependence introduces systematic bias into multi-agent settings. \rev{Language may also change the effect of the stakes, not only the level of cooperation: human responses to stake size differ across societies~\cite{bosman2017emotion,cameron1999raising}, and the cultural values that language models express shift with the prompt~\cite{tao2024cultural}, so a change of scale may be read differently in each language.}


This dependence is an instance of a more general problem: measurements of language-model behaviour are often sensitive to aspects of the elicitation procedure that are not themselves the intended object of measurement. Accuracy can vary with the ordering of in-context examples~\cite{lu2022fantastically}, and with formatting choices that preserve the semantic content of the prompt, by margins large enough to obscure the differences between the models under comparison~\cite{sclar2024quantifying}. Persona instructions are similarly consequential: they alter model outputs~\cite{deshpande2023toxicity}, affect measured capability~\cite{salewski2023context}, and, when specified through a demographic backstory, can induce response distributions resembling those of the corresponding human subgroup~\cite{argyle2023out}. They therefore do not establish a well-defined null condition against which the effect of prompting can be assessed.


The choice of units (of measurement)  is rarely presented as a design decision, because in the standard theory it is not one. \rev{When the entries of a payoff matrix are read as von Neumann-Morgenstern utilities, the preferences they represent are unchanged by any positive affine transformation}~\cite{von2007theory,luce1957games}. Multiplying every entry in the payoff matrix by ten changes no preference ordering, no best response, no equilibrium, and no \rev{orbit} of the replicator dynamics\rev{; it rescales only time, so convergence speeds up or slows down along the same path}. The result is not a new game with larger payoffs, but the same game written in different units.  The theory of the iterated dilemma inherits this invariance: its predictions depend on the ordering of the  payoff-matrix entries  and on ratios among them, rather than on their absolute magnitudes~\cite{axelrod1981evolution,nowak1992tit,nowak1993strategy,pressDyson2012}. For this reason, the modelling literature often normalises the payoff scale away~\cite{szabo2007evolutionary,perc2010coevolutionary,wang2015universal}. \rev{Recent} theory strengthens this point: in models in which social learners update on the payoff they most recently received, the condition for conditional cooperation to persist under strong selection contains neither the benefit nor the cost of cooperation, and depends only on the ordinal structure of the payoffs~\cite{glynatsi2024evolution}.

The invariance is a property of the theory, not of human decision-making. Human valuation is neither linear in magnitude nor anchored at zero~\cite{kahneman1979prospect}, human choice reverses between logically equivalent descriptions of one problem~\cite{tversky1981framing}, and human cooperation in the dilemma itself responds to the size of the stakes~\cite{krockow2016cooperation,list2006friend}. \rev{These effects are usually modest and depend on the task and on stakes that are large relative to income~\cite{camerer1999effects,kocher2008does,engel2011dictator,larney2019stake,cameron1999raising,andersen2011stakes,holt2002risk}.} That gap has been argued to be structural rather than incidental, with the linguistic description of an action entering the utility function directly and a shift proposed from outcome-based to language-based preferences~\cite{capraro2024language}.

Two things are established, and they have not been put together. The first is that what a language model does depends on features of the prompt that carry no task content: the order of the examples, the formatting, the persona, the language. The second is that multiplying every payoff-matrix entry by a positive constant carries no task content in the strictest sense available, because the matrix it produces is  the same game. What is not known is whether the strategic behaviour of a language model respects that invariance. To our knowledge, no prior study has formulated positive payoff rescaling as a theorem-backed zero-effect invariance test across a broad multi-model, multilingual grid. Every manipulation in the sensitivity literature could, in principle, induce legitimate behavioural changes. Consequently, observed changes under these manipulations are difficult to interpret, since the expected direction or magnitude of the effect is typically unspecified. Such changes may therefore be treated either as methodological noise or as empirical findings, but only relative to an implicit and under-theorised baseline. Payoff rescaling is distinctive because its theoretical prediction is explicit: behaviour should remain unchanged. A measured departure from this prediction therefore has a clearly defined null reference, rather than only a contrast with another experimental condition.

% What follows if the invariance fails is a question about what a published
% cooperation rate means. Machine behaviour proposes that autonomous systems be
% studied empirically, with the instruments of the behavioural sciences, because
% what such a system does is not deducible from how it was
% built~\cite{rahwan2019machine}. Those instruments arrive with a known
% difficulty: a behavioural manipulation almost never carries a predicted effect,
% so a measured change has to be considered against a comparison condition the
% experimenter chose. The invariance classes of decision theory are the
% exception, supplying a manipulation whose predicted effect is a stated number,
% and the number is zero. The consequences run through the strands of machine
% behaviour this study touches. The persona line is the one place where the
% experimenter addresses the agent about its social role, and whether it takes
% hold turns out to depend on the units the matrix was shown in. Models of
% mixed human-machine populations carry the cooperativeness of the artificial agent as
% a parameter and ask where the population
% settles~\cite{santos2018social,han2011intention,han2013good,perc2017statistical,paiva2018engineering};
% a rate partly determined by the experimenter's units is not the quantity those
% models need.



Taken together, these considerations motivate three questions about the effects of payoff rescaling:
\begin{enumerate}[(1)]
\item Is strategic behaviour invariant to   rescaling of payoffs, and how large are any deviations from invariance?
\item Does the result depend on the language of the prompt and on the persona assigned to the model?

\item Does payoff rescaling affect only how often the model  cooperates, or does it also alter the underlying strategy it adopts?

\end{enumerate}
We systematically address these questions on a complete and exactly balanced grid of 10{,}000 configurations (dyads) (Section \ref{sec:methodology}), and compare the responses against an evolutionary game baseline computed on the same grid.

%Here, we turn a theorem into an instrument. 
Our main contributions  are: \textbf{(i)} A payoff-rescaling protocol whose predicted behavioural effect is exactly zero, demonstrating that five widely adopted commercial language models systematically violate scale invariance (Section \ref{sec:payoff_magnitude}), with cooperation rates shifting substantially across payoff scales. \textbf{(ii)} A per-model analysis showing that scale departures diverge sharply across models to the point of anti-correlation. Consequently, aggregate analyses produce a synthetic average that is representative of no individual model. Variance is dominated by model identity and model-by-scale interactions; while the common scale effect in the five-model decomposition is largely driven by Grok, the aggregate dyad-level association persists in leave-one-model-out analyses. \textbf{(iii)} Cross-lingual and structural evidence demonstrating that payoff-scale sensitivity is strongly language-dependent, modulates persona-induced behaviour at sub-unit stakes, and alters the strategic profiles exhibited by the models. \textbf{(iv)} An analytical evolutionary game baseline where evolutionary outcomes diverge from the aggregate empirical patterns, and a transition analysis demonstrating  that own-action  conditioning generally outweighs partner-action conditioning (Section \ref{sec:conditioning}, appendix~D.7).



% ===========================================================================
\section{Material and methods}\label{sec:methodology}

\subsection{Framework}\label{sec:background}

FAIRGAME, the Framework for AI Agents Bias Recognition using Game Theory~\cite{buscemi2025fairgame}, supplies the computational infrastructure for this study. It supports systematic and reproducible evaluation of language models through controlled game-theoretic experiments: an experimental condition is declared in a configuration file that fixes the payoff structure, the horizon, the model backend and the prompt language; at run time the framework combines that declaration with language-specific prompt templates, simulates the repeated normal-form game, and logs the round-by-round trajectory. We extend it in three ways. First, with a payoff-scaling factor applied to the prisoner's dilemma, which is the manipulation this paper examines (Section \ref{sec:gameImportance}). Second,  we apply a completeness safeguard that prevents the release of an analysis table unless the experimental grid is exactly balanced and the recorded payoffs for each game correctly recover the payoff scale assigned to its corresponding cell (i.e. a  combination of experimental conditions, such as model, language, persona, and payoff scale). Third, we include a strategy read-out that identifies which canonical memory-one rule each trajectory most closely resembles, while also reporting the provenance of that classification (Section \ref{sec:readout}).

The design is a complete factorial crossing of five language models, ten payoff scales, five prompt languages, four persona pairings and ten replicates. The unit of collection is the dyad, one ten-round game between two agents of the same model, contributing two agent-games, an agent-game being the ten decisions one agent takes in one game. Each of the 250 model-by-language-by-scale cells holds exactly 40 dyads, the four persona pairings by ten replicates; pooling the five languages gives 200 dyads and 400 agent-games in each of the 50 model-by-scale cells. The main-text corpus is 10{,}000 dyads, 20{,}000 agent-games and 200{,}000 decisions, and is the grid every main-text figure and table is computed on. Balance is enforced rather than assumed (appendix~A.1).

\subsection{The payoff-scaled prisoner's dilemma and its invariance}\label{sec:gameImportance}

Both agents receive the prisoner's-dilemma prompt of the FAIRGAME protocol~\cite{buscemi2025fairgame} in its no-communication, announced-horizon condition. The scenario casts the strategic interaction as a criminal interrogation, states the four payoff outcomes as penalties, names the two available actions Option~A and Option~B, carries the running history of prior moves and realised penalties, announces a fixed horizon of $T = 10$ rounds, and instructs each agent to minimise its cumulative penalty.

The row player's baseline payoff matrix for the Prisoner's Dilemma is
\begin{center}
\small
\begin{tabular}{c|cc}
 & Option~A & Option~B \\
\hline
Option~A & $(6, 6)$ & $(0, 10)$ \\
Option~B & $(10, 0)$ & $(2, 2)$
\end{tabular}
\end{center}
where the first and second entries in each cell denote the penalty payoffs to the row and column player, respectively. In standard game-theoretic notation, this matrix assigns the four canonical outcomes:
\[
T = 0, \quad R = 2, \quad P = 6, \quad S = 10,
\]
representing the temptation ($T$), reward ($R$), punishment ($P$), and sucker ($S$) penalty payoffs. The penalty ordering
\[
T < R < P < S \quad \text{with} \quad 2R < T + S
\]
is the exact isomorphic image of the familiar canonical Prisoner's Dilemma condition $T_{\text{util}} > R_{\text{util}} > P_{\text{util}} > S_{\text{util}}$ and $2R_{\text{util}} > T_{\text{util}} + S_{\text{util}}$ defined on utilities $u = -\text{penalty}$:
\[
T_{\text{util}} = 0 > R_{\text{util}} = -2 > P_{\text{util}} = -6 > S_{\text{util}} = -10.
\]
Because lower penalties correspond to higher utilities, Option~A strictly dominates Option~B in every column (facing Option~A, choosing A yields penalty 6 versus 10; facing Option~B, choosing A yields penalty 0 versus 2). Option~A is therefore the defecting action, while mutual Option~B yields penalty 2 versus penalty 6 for mutual Option~A, making Option~B the cooperative action; throughout this paper, cooperation denotes choosing Option~B.

The realised score-action relationships reproduce the intended dilemma structure in the logs: own cooperation correlates with own utility at Spearman $\rho = -0.015$, while the opponent's cooperation correlates at $\rho = +0.916$, and dyad welfare correlates with joint cooperation at $\rho = +0.985$. These checks validate the payoff mapping and action polarity, not model comprehension of the dilemma (appendix~A.1).

To manipulate the stakes of the game without altering its strategic structure, best-reply correspondences, or Nash equilibria, we introduce a positive scalar parameter $\lam > 0$ and multiply all four baseline cells by $\lam$ (figure~\ref{fig:overview}). In our experiments, we evaluate ten values spanning five orders of magnitude:
\[
\lam \in \{0.01,\ 0.10,\ 0.25,\ 0.50,\ 1.0,\ 2.0,\ 5.0,\ 10.0,\ 100.0,\ 1000.0\},
\]
corresponding to attenuated ($\lam < 1.0$), baseline ($\lam = 1.0$), and amplified ($\lam > 1.0$) payoff regimes. For example, at the sub-unit scale $\lam = 0.10$, the row player's payoff matrix becomes
\begin{center}
\small
\begin{tabular}{c|cc}
 & Option~A & Option~B \\
\hline
Option~A & $(0.6, 0.6)$ & $(0, 1.0)$ \\
Option~B & $(1.0, 0)$ & $(0.2, 0.2)$
\end{tabular}
\end{center}
where every  payoff entry is at most 1.0. At the amplified scale $\lam = 10.0$, it becomes
\begin{center}
\small
\begin{tabular}{c|cc}
 & Option~A & Option~B \\
\hline
Option~A & $(60, 60)$ & $(0, 100)$ \\
Option~B & $(100, 0)$ & $(20, 20)$
\end{tabular}
\end{center}
with the ordering $T < R < P < S$ strictly preserved in all cases. The complete matrix for all ten scales is provided in appendix~A.3.

Under classical decision theory and game theory, multiplying utilities by $\lam > 0$ constitutes a positive affine transformation of a von Neumann-Morgenstern utility function~\cite{von2007theory,luce1957games}:
\[
u'(a) = \lam u(a) + c \quad (\lam > 0,\ c = 0).
\]
Such transformations preserve all preference orderings over lotteries, best replies, Nash equilibria, and deterministic replicator dynamics orbits. Consequently, the normative theoretical prediction for the effect of $\lam$ on rational agents is exactly zero. \rev{This is a result about the game. It predicts unchanged behaviour only if the stated penalties are the agent's whole utility and its choices do not depend on their numerical form; the criminal-interrogation narrative and the persona instruction do not scale with $\lam$ and may add considerations beyond the penalties. The null is therefore a prediction for an agent that meets these assumptions.} This construction isolates the effect of payoff magnitude from structural incentives: any systematic departure from invariance constitutes an empirical discovery about machine behaviour rather than a rational adaptation to changing game-theoretic payoffs.

\begin{figure}[!htbp]
\centering
\includegraphics[width=\linewidth]{figures/prisoners_dilemma_overview.pdf}
\caption{\textbf{Overview of the repeated prisoner's-dilemma experiment.} Two language-model agents from the same provider repeatedly play the stage game across ten payoff scales, with penalties to minimise. The main-text corpus covers five widely adopted commercial models, five prompt languages, four persona pairings and 10{,}000 dyads (20{,}000 agent-games and 200{,}000 decisions).}
\label{fig:overview}
\end{figure}

\subsection{Models, prompt languages, personas and protocol}\label{sec:protocol}

Five widely adopted commercial models were run, each against itself; two different models were never paired. They were Claude Haiku 4.5, GPT-5.4 Nano, Gemini 3.5 Flash-Lite, Qwen3 235B-A22B and Grok 4.20. Their model identifiers are listed in appendix~A.2.

The prompt languages are English, French, Vietnamese, Chinese and Arabic, in templates byte-identical to the FAIRGAME resource files and checked against them before every collection run. Artefacts of those files are reproduced rather than repaired, the consequential one being that the Arabic and Chinese versions state the four cells as penalties but phrase the goal sentence as maximising a reward, so we treat the five languages as five distinct stimuli rather than five renderings of one (appendix~A.4). Each agent also receives one line naming its persona, in the language of the prompt: in English, \emph{You are cooperative.} or \emph{You are selfish.} The two are crossed in all four ordered pairings, ten replicates each. The opponent's persona is never disclosed, so a pairing reaches a player only through realised behaviour and not before round 2, and we treat it as a design factor and a regression covariate rather than as information the agent held.

Every game runs ten rounds with simultaneous moves: both agents see the same history, rounds 1 to $r-1$, and the round is scored once both have replied. Each game is repeated at all ten payoff scales with the same persona pairing and the same requested sampling seeds (appendix~A.2). Decisions are elicited as free text and parsed by the protocol's own matcher. A reply it cannot resolve is regenerated, up to two further attempts, and if all three fail the decision is recorded as Option~A, which under this framing is defection: the fallback can only push a measured cooperation rate down, never up, and it cannot manufacture the cooperative behaviour reported here. Each collection run asserts that its fallback rate is at most 0.02\rev{; the per-game tables do not record the number of attempts, so the rates cannot be broken down by language and scale} (appendix~D.4). \rev{The cost of the collection is reported in appendix~A.2.}

\subsection{Strategy read-out}\label{sec:readout}

While the framework provides complete gameplay trajectories and descriptive metrics such as cooperation rates, these capture what agents do rather than the decision rule behind it. Following prior work on inferring canonical strategies from noisy repeated-game data~\cite{han2013intention,Han2012ALIFEjournal,diStefanoIntention2023}, we read every agent-game against the four canonical rules AllC, AllD, Tit-for-Tat and Win-Stay-Lose-Shift. The read-out reports which of the four a ten-round trajectory most \emph{resembles}. It is not an identification of the strategy a model is running, and no claim in this paper treats it as one.

All four rules are memory-one, so exact matching is a lookup with no learning in it, and it returns a \emph{set} rather than a label, because over ten rounds several rules routinely prescribe the same moves. The three provenances, the agent-games matched exactly by one rule (deduced), by several at once (ambiguous) and by none (unmatched), are recorded separately and never merged, and the same matcher yields a continuous quantity needing no labels at all, the number of rounds violating the nearest rule. Where the rules leave the answer open a single-layer LSTM~\cite{lstm} arbitrates. Trained on a synthetic corpus alone, which contains no language-model output, it reaches 98.5\% accuracy in distribution and 93.6\% on a held-out execution-noise condition (appendix~A.6). It ranks only within the set the rule base returns, so it can never overturn a deduction; it decides alone only where no rule fits, and its share of the labels is very unevenly spread across the four rules, which bounds how the composition may be read (Section \ref{sec:llm_strategies}).

\subsection{Statistical analysis}\label{sec:stats}

The cooperation rate of an agent-game is the share of its ten decisions that are Option~B; a cell mean averages the agent-games of a cell, and round-level analyses use the binary decision itself. Because the two agents of a dyad play each other and the ten rounds inside an agent-game are strongly autocorrelated, primary inferential intervals were computed by resampling whole dyads with replacement, 1{,}500 draws, percentile method~\cite{efron1993introduction}, and primary per-model and trend regressions clustered their standard errors on the dyad. \rev{Across scales, the ten dyads that share a model, a language and a game identifier, and so a persona pairing and the requested seeds, form a matched ladder, 200 per model. Appendix~D.6 repeats the headline inferences with the ladder as the unit; no conclusion changes.} The supplementary scale-by-language robustness model instead clusters on the documented language-by-game-ID blocks to preserve the matched ten-scale ladders. Descriptive figure intervals use whole-dyad resampling. Because the range over the ten scales is a maximum minus a minimum, and so positive by construction, its bootstrap interval quantifies the precision of the range and not its distance from zero; the permutation test supplies the test. Because the null is a theorem, the matching test destroys the association with the payoff scale and leaves the rest of the design as observed: the scale label is shuffled among whole dyads, the unit to which the scale was assigned, and the range of the ten scale means recomputed, 2{,}000 draws, so $0.0005$ is the smallest attainable $p$ value.

For each model we fitted a logistic regression of the round-level decision on the payoff scale as a ten-level factor, with prompt language, own persona and opponent persona as covariates, $\lam = 1$ as reference and standard errors clustered on the dyad, and tested the nine scale contrasts jointly by a Wald test (table~\ref{tab:effect}). The variance decomposition is a type-II analysis of variance on the 250 cell means of the five models reported here, with terms for model, language and payoff scale and their three two-way interactions and no three-way term, so the residual on 144 degrees of freedom is the omitted three-way interaction; the supplementary decomposition is in appendix~B.3. The persona result contrasts the two sub-unit scales with the remaining eight as a difference of differences with an interval from the same dyad bootstrap; where a trend across the whole grid is wanted instead, $\log_{10}\lam$ enters as a continuous predictor. Analyses were run in Python; the package versions and the seeds are recorded with the archived pipeline in appendix~E.


% ===========================================================================
\section{Results}\label{sec:results}

The results below are computed on the 200{,}000 recorded decisions in the five-model main-text corpus.

\subsection{Cooperation varies under payoff rescaling }\label{sec:payoff_magnitude}

Under normative decision theory and classical game theory, multiplying every payoff in a normal-form matrix by a positive constant is an axiomatically neutral transformation. It preserves all preference rankings over lotteries, best-reply correspondences, and Nash equilibria: an agent playing the game should not alter its choices whether payoffs are denominated in cents, dollars, or thousands of pounds.

Yet, across five orders of payoff magnitude, every language model evaluated violates this scale invariance (see figure~\ref{fig:landscape}a,c, table~\ref{tab:effect}). The shift in strategic behaviour is substantial: across the ten-scale ladder, cooperation varies markedly for every model, representing a clear departure from the scale-invariance null (see figure~\ref{fig:landscape}b and table~\ref{tab:effect}). What classical game theory treats as an inert change of units produces a fundamental reordering of multi-agent behaviour.

Notably, the models do not share a universal behavioural pattern: each possesses a distinct machine character that responds idiosyncratically to the scale of the stakes. Grok exhibits a \rev{strong behavioural} shift across payoff scales. At micro-stakes, it predominantly  defects; at standard stakes, however, it shifts towards high levels of cooperation, selecting the socially optimal action in the vast majority of interactions. GPT climbs sharply from a hesitant baseline before saturating at a stable cooperative tendency. Claude’s cooperation rate decreases  at intermediate sub-unit stakes before recovering at larger payoff scales. Gemini maintains an increased  cooperative tendency throughout, while Qwen becomes slightly less cooperative as the stakes rise (most of the \rev{change} occurs between $\lam = 0.1$ and $\lam = 1$). Because each model follows a distinct behavioural trajectory, machine cooperation cannot be characterised by a single universal behavioural pattern.

The largest departures are concentrated near the low-payoff end of the grid (table~\ref{tab:effect}). Fractional payoff multipliers often reduce cooperation compared with intermediate scales, whereas changes from unit stakes to large stakes have no consistent pattern across models. This result should be interpreted descriptively. The present design does not establish a discrete cognitive threshold, distinguish numerical magnitude from every feature of notation, or identify the mechanism behind the response.

\begin{figure}[!htbp]
\centering
\includegraphics[width=\linewidth]{figures/f_landscape.pdf}
\caption{\textbf{Cooperation rates as a function of the payoff scale across five models.} (a) Model-by-scale cooperation matrix: mean cooperation for each of the five models at each of the ten payoff scales $\lam \in [0.01, 1000]$, with the value shown in each box. (b) Within-model cooperation range (maximum minus minimum rate across the ten scales) compared against the pooled five-model range (0.154, vertical dashed line). \rev{The pooled range is narrower than the within-model range of three of the five models.} Confidence intervals are 95\% percentile intervals from 1{,}500 whole-dyad bootstrap resamples. (c) Small multiples per model: the focal model is drawn with its designated colour and marker, with the other four models shown behind in grey for context. Shaded bands mark the sub-unit regime ($\lam \leq 0.1$). Each model-by-scale cell pools 200 dyads (400 agent-games), 20{,}000 agent-games in all.}
\label{fig:landscape}
\end{figure}

% ===== begin tables_auto.tex (inlined) =====
% Generated by Analysis/scaling/s06_tables.py - do not edit by hand.
\begin{table}[t]
\centering
\caption{\textbf{Cooperation over the ten payoff scales, by model.} Minimum, maximum and range, with a 95\% confidence interval on the range from a bootstrap over whole dyads and a permutation test on the payoff-scale label ($0.0005$ is the floor of 2{,}000 draws). The last two columns test the nine scale contrasts jointly in a per-model logistic regression of the round-level decision, clustered on the dyad. Each entry pools 200 dyads per scale; ranges are computed before rounding. Full specification in the electronic supplementary material.}
\label{tab:effect}
\small
\begin{tabular}{lccccccc}
\toprule
& \multicolumn{2}{c}{cooperation} & & & & \multicolumn{2}{c}{Wald test} \\
\cmidrule(lr){2-3} \cmidrule(lr){7-8}
Model & min & max & range & 95\% CI & $p_{\text{perm}}$ & $\chi^2_9$ & $p$ \\
\midrule
Claude Haiku 4.5 & 0.371 & 0.452 & 0.081 & [0.062,\,0.135] & $0.0025$ & 32.3 & $<0.001$ \\
GPT-5.4 Nano & 0.438 & 0.680 & 0.243 & [0.200,\,0.290] & $0.0005$ & 164.1 & $<0.001$ \\
Gemini 3.5 Flash-Lite & 0.682 & 0.872 & 0.190 & [0.165,\,0.240] & $0.0005$ & 129.0 & $<0.001$ \\
Qwen3 235B-A22B & 0.399 & 0.494 & 0.095 & [0.068,\,0.158] & $0.0005$ & 96.9 & $<0.001$ \\
Grok 4.20 & 0.274 & 0.910 & 0.636 & [0.588,\,0.686] & $0.0005$ & 1227.3 & $<0.001$ \\
\bottomrule
\end{tabular}
\end{table}

\begin{table}[t]
\centering
\caption{\textbf{Persona effect below and above the sub-unit boundary, by model.} Cooperation under a cooperative persona minus cooperation under a selfish one, at the two scales where every shown payoff is at most 1 ($\lam \leq 0.1$) and at the eight where some payoff exceeds it. Intervals are 95\% confidence intervals from the dyad bootstrap. A sign flip needs both regime intervals to exclude zero on opposite sides, so a supra-unit interval straddling zero means the effect is abolished rather than reversed.}
\label{tab:persona}
\footnotesize
\setlength{\tabcolsep}{3pt}
\begin{tabular}{lccccc}
\toprule
& \multicolumn{2}{c}{persona effect, with 95\% CI} & & & \\
\cmidrule(lr){2-3}
Model & $\lam \leq 0.1$ & $\lam \geq 0.25$ & shift & 95\% CI & sign flip \\
\midrule
Claude Haiku 4.5 & $-0.222$ $[-0.258,\,-0.186]$ & $+0.207$ $[+0.189,\,+0.226]$ & 0.429 & [0.389,\,0.471] & yes \\
GPT-5.4 Nano & $-0.135$ $[-0.179,\,-0.091]$ & $+0.018$ $[-0.001,\,+0.037]$ & 0.153 & [0.105,\,0.200] & no \\
Gemini 3.5 Flash-Lite & $-0.456$ $[-0.507,\,-0.406]$ & $+0.224$ $[+0.207,\,+0.242]$ & 0.680 & [0.625,\,0.736] & yes \\
Qwen3 235B-A22B & $-0.911$ $[-0.931,\,-0.893]$ & $-0.679$ $[-0.694,\,-0.664]$ & 0.233 & [0.210,\,0.258] & no \\
Grok 4.20 & $-0.441$ $[-0.485,\,-0.396]$ & $-0.303$ $[-0.326,\,-0.279]$ & 0.138 & [0.090,\,0.185] & no \\
\bottomrule
\end{tabular}
\end{table}

\begin{table}[t]
\centering
\caption{\textbf{Rule agreement and rule distance against the payoff scale, by model.} Left, the share of agent-games matched exactly by one canonical rule, below and above the sub-unit boundary, with the coefficient of a logistic regression of that indicator on $\log_{10}\lam$; right, the mean number of rounds violating the nearest canonical rule at the smallest and largest scale, with the same trend by least squares. Both $\beta$ are per decade of $\lam$ and their signs differ between models. Both cluster on the dyad, and neither block uses the learned classifier.}
\label{tab:strategy}
\small
\setlength{\tabcolsep}{4pt}
\begin{tabular}{lcccccccc}
\toprule
& \multicolumn{4}{c}{exactly one rule fits (\%)}
& \multicolumn{4}{c}{rounds violating the nearest rule} \\
\cmidrule(lr){2-5} \cmidrule(lr){6-9}
Model & $\lam \leq 0.1$ & $\lam \geq 0.25$ & $\beta$ & $p$
      & $\lam{=}0.01$ & $\lam{=}10^3$ & $\beta$ & $p$ \\
\midrule
Claude Haiku 4.5 & 17.9 & 3.5 & $-0.545$ & $<0.001$ & 1.70 & 2.02 & $+0.073$ & $<0.001$ \\
GPT-5.4 Nano & 9.8 & 7.6 & $-0.038$ & $0.441$ & 2.02 & 2.19 & $+0.038$ & $0.052$ \\
Gemini 3.5 Flash-Lite & 32.2 & 15.6 & $-0.242$ & $<0.001$ & 0.48 & 1.36 & $+0.194$ & $<0.001$ \\
Qwen3 235B-A22B & 72.1 & 53.6 & $-0.267$ & $<0.001$ & 0.27 & 1.14 & $+0.202$ & $<0.001$ \\
Grok 4.20 & 44.6 & 17.1 & $-0.301$ & $<0.001$ & 0.92 & 0.37 & $-0.118$ & $<0.001$ \\
\bottomrule
\end{tabular}
\end{table}
% ===== end tables_auto.tex =====

\subsection{Effect sizes vary across prompt languages}\label{sec:language}

Prompt language is not a neutral transmission channel for the game. Both the baseline level of cooperation and its sensitivity to the stakes diverge substantially across prompt languages. Figure~\ref{fig:language} shows the pooled language profiles, while the model-by-language levels and ranges appear in electronic supplementary table~B.3.

We note a caveat related to the comparison of levels rather than the scale effect itself. The Arabic and Chinese templates differ from the other prompt stimuli in their objective wording. They state the four cells as penalties while phrasing the goal sentence as maximising a reward (Section \ref{sec:protocol}). Because that wording is fixed within each language, the within-language comparisons still establish sensitivity to payoff rescaling for the prompt as actually presented. However, differences in the magnitude or shape of scale sensitivity across languages cannot be attributed to language alone, because language is partially confounded with template wording. We therefore interpret the scale-by-language interaction as a prompt-stimulus interaction, not as evidence for a linguistic or cultural mechanism.

The central finding is that scale sensitivity itself depends on the prompt language. The max-minus-min range over the ten scales differs between languages by a factor of about two within every model, and by an order of magnitude across model-language cells (electronic supplementary table~B.3). Figure~\ref{fig:language}b provides the complementary supra-unit minus sub-unit contrast, and figure~\ref{fig:language}c shows that language differences remain after removing their mean levels. Across the main-text corpus, the observed sensitivity ranges from near-invariance to sweeping behavioural transformations spanning almost the entire probability scale.

Crucially, these linguistic biases do not point in a uniform direction across models: the most cooperative language differs from model to model, whereas English is the most scale-sensitive language for four of the five models. When multilingual evaluations pool across models, these opposing model-specific effects are attenuated and can obscure substantial model-by-language heterogeneity. In reality, machine behaviour depends on the linguistic stimulus used in the prompt. Within this design, scale sensitivity is a joint property of the model and prompt language rather than a model-independent invariant. It does not identify a cultural mechanism.

\begin{figure}[!htbp]
\centering
\includegraphics[width=\linewidth]{figures/f_language.pdf}
\caption{\textbf{Cross-lingual scale sensitivity.} (a) Cooperation pooled across the five main-text models for each prompt language and payoff scale; the dashed boundary separates the sub-unit and supra-unit regimes. (b) Language-specific difference in cooperation between supra-unit and sub-unit regimes; intervals show the pooled estimate and small markers the five models separately. (c) Cooperation curves after subtracting each language's mean, showing that differences in scale-response shape remain after level differences are removed.}
\label{fig:language}
\end{figure}

\subsection{Model heterogeneity and the effects of aggregation}\label{sec:heterogeneity}

Averaging across models conceals the real phenomenon through mutual cancellation. The pooled response curve severely attenuates the true magnitude of scale sensitivity, shrinking the typical within-model movement substantially (figure~\ref{fig:landscape}b). Because pairwise correlations between model curves are predominantly negative (figure~\ref{fig:landscape}c), curves that move in opposite directions flatten each other out when summed. A pooled curve describes an artificial average agent that no real model resembles.

Underlying this cancellation is a fundamental divide in policy geometry: cooperation in these models is not a single Gaussian distribution, but a mixture over distinct policy regimes (appendix~B.6). Across the corpus, unconditional boundary play is strongly overrepresented, with extreme atoms carrying far more weight than expected under uniform play. Crucially, the commitment to boundary play varies dramatically across models: Claude operates as an interior, mixed-policy player that rarely locks into extreme  behavioural patterns, whereas Qwen and Grok function as pure-policy switchers, oscillating between unconditional defection and unconditional cooperation. Pooling across models conflates how often a model chooses a pure strategy with which pure strategy it selects.

A formal variance decomposition confirms that model individuality dominates the game (table~B.4, appendix~B.3). Model identity and the model-by-language interaction account for the majority of the variance across cells. While the main effect of payoff scale is statistically detectable, its interaction with model identity carries three times as much variance, demonstrating that the behavioural impact of stakes is model-specific rather than shared universally across model endpoints.

In the five-model cell-mean variance decomposition, the shared payoff-scale main effect is largely carried by Grok. Removing Grok renders that particular main-effect test non-significant, even though every remaining model continues to move (appendix~B.3). A complementary dyad-level leave-one-model-out test asks a different question, whether any pooled scale association remains. Every five-model subset retains a detectable association, including the subset without Grok (appendix~D). The two analyses are therefore complementary: one tests a shared main effect after accounting for heterogeneity, whereas the other tests any residual pooled association.

\subsection{Payoff scale modulates persona-induced cooperative behaviour}\label{sec:persona}

System prompts designed to steer agent behaviour do not operate independently of game stakes. The behavioural difference between cooperative and selfish persona instructions changes dramatically across the payoff ladder (figure~\ref{fig:persona}, table~\ref{tab:persona}). This modulation differs between the two sub-unit scales and the eight larger scales, where shown payoffs include values above one (Section \ref{sec:gameImportance}). Across all five models, the persona effect changes across these two regimes.

In several models, the persona instruction reverses its behavioural effect  below the unit threshold. In Gemini and Claude, the persona effect reverses sign: under fractional stakes, both models exhibit higher cooperation when explicitly instructed to behave selfishly than when instructed to behave cooperatively. In GPT, the selfish persona also raises cooperation at fractional stakes, but above them the persona effect disappears instead of reversing, while the remaining models maintain their pattern. This result is a boundary-aligned descriptive effect. It does not identify why the prompt instruction changes direction at these scales.

The main grid does not separate numerical magnitude from all aspects of numeral formatting. We therefore treat this regime split as descriptive and do not attribute it to a specific representational mechanism.

A separate class of models inverts the persona instruction across the entire ladder (table~\ref{tab:persona}). Qwen and Grok cooperate substantially more under the selfish persona than under the cooperative persona at every evaluated scale. These models show an opposite response to the persona wording. The data do not establish why they do so. This result underscores that role prompts cannot be assumed to steer multi-agent interactions in the intended direction without empirical verification across payoff contexts (appendix~B.4).

\begin{figure}[!htbp]
\centering
\includegraphics[width=\linewidth]{figures/f_persona.pdf}
\caption{\textbf{Persona effect and sub-unit inversion against the payoff scale.} The persona effect is the cooperation rate under the cooperative persona minus that under the selfish persona. (a) Difference in cooperation rate across scales $\lam$, with the sub-unit regime ($\lam \le 0.1$) shaded where all five models cooperate more under the selfish prompt. (b) Regime means comparing sub-unit ($\lam \le 0.1$, open circles) with supra-unit ($\lam \ge 0.25$, filled circles) regimes (table~\ref{tab:persona}).}
\label{fig:persona}
\end{figure}

\subsection{ Payoff scale alters strategic composition, not merely cooperation rates}\label{sec:llm_strategies}

Two models displaying identical cooperation rates can be executing fundamentally different behavioural strategies/rules. To understand the behavioural rules underlying raw cooperation rates, we read every agent trajectory against the four canonical memory-one strategies of evolutionary game theory (Section \ref{sec:readout}). What this read-out captures is behavioural resemblance rather than exact algorithmic identity: across the corpus, the majority of games match no textbook rule cleanly (table~\ref{tab:strategy}, appendix~B.5).

As payoff stakes increase, play moves farther from the four canonical rules in most models. Across the scale ladder, the proportion of games adhering strictly to a single canonical rule declines substantially, while unmatched play accumulates (table~\ref{tab:strategy}). This loss of canonical-rule conformity is pronounced in Claude, Grok, Qwen, and Gemini, while GPT maintains steady adherence. The read-out does not establish what rule generates unmatched play; it only shows that it is farther from the four canonical rules.

A label-free continuous metric confirms this behavioural drift: the number of rounds departing from the nearest canonical rule increases significantly with scale across most models (table~\ref{tab:strategy}). As stakes multiply, trajectories depart more often from the textbook templates. Only Grok reverses this trend, moving towards extreme pure-policy play at high stakes.

The qualitative composition of strategies also shifts dynamically, driving divergent realignments across models (figure~\ref{fig:mix}). In GPT, expanding stakes prompts a steady migration from unconditional defection toward unconditional cooperation. In Gemini, the opposite shift occurs: unconditional cooperation yields ground to more Tit-for-Tat-like and Win-Stay-Lose-Shift-like attributions. Grok undergoes an outright reversal, swapping rampant defection at micro-stakes for overwhelming cooperation at high stakes.

Critically, these strategy shifts reflect qualitative resemblances rather than exact implementations \rev{of the canonical strategies}. Provenance auditing demonstrates that while unconditional cooperation and defection are often matched exactly, nearly all attributions of Tit-for-Tat and Win-Stay-Lose-Shift represent best-fit approximations to noisy, unmatched play (figure~\ref{fig:mix}c, appendix~B.5). When models appear to increase their Tit-for-Tat share, they are not adopting the classic reciprocal algorithm; they are generating complex, imperfect trajectories whose overall statistical contours happen to lean in that direction.

\begin{figure}[!htbp]
\centering
\includegraphics[width=\linewidth]{figures/f_strategy_mix.pdf}
\caption{\textbf{Strategy attribution across payoff scales, composition shift, and provenance breakdown.} (a) Stacked area share of canonical memory-one strategy attributions (AllC, AllD, TFT, WSLS) across the ten payoff scales for each of the five models. (b) Net strategy share changes between extreme scales ($\lam = 10^{3}$ minus $\lam = 0.01$) across the four canonical rules with total variation distance $D_{\mathrm{TV}}$ in parentheses. (c) Strategy read-out provenance: exactly deduced (single canonical match), ambiguous (multiple matching rules), and unmatched (LSTM attribution on trajectories consistent with no standard rule). Over the corpus, $94.6\%$ of Tit-for-Tat and $97.1\%$ of Win-Stay-Lose-Shift labels are attributions on unmatched play, establishing that labels indicate resemblance rather than faithful execution.}
\label{fig:mix}
\end{figure}

\subsection{Where payoff scale  acts: the evolutionary baseline}\label{sec:egt_pointer}

While payoff scale is theoretically inert in normal-form games, evolutionary game theory provides a foundational framework where stakes legitimately matter, offering a rigorous theoretical benchmark for language model behaviour.

In deterministic replicator dynamics, scale is completely neutral. Multiplying the payoff matrix by a positive constant simply scales the velocity of the vector field: orbits follow identical geometric trajectories and all rest points remain invariant (appendix~C.2, figure~C.1).

In finite populations with stochastic imitation \citep{nowak2004emergence}, however, payoff scale acts directly through the selection intensity. Under pairwise Fermi imitation, scaling payoffs amplifies the competitive advantage of defection over cooperation (appendix~C).  \rev{With $\beta$ fixed, payoffs enter the Fermi rule only through $\beta\lam$, so raising $\lam$ is equivalent to raising the intensity of selection.} At micro-stakes, selection is weak and all four canonical rules coexist near genetic drift. As stakes expand, selection pressure intensifies: unconditional defectors strongly invade cooperative strategies until AllD achieves near-total fixation (figure~\ref{fig:egtllm}a). Combining evolutionary dynamics with language-model evaluations provides a principled lens for multi-agent governance~\cite{balabanova2025media,buscemi2025llms,xie2026evolutionary,han2026social}.

At the pooled five-model level, the AllD label share moves in the opposite direction to the finite-population evolutionary baseline as $\lam$ increases (figure~\ref{fig:egtllm}). This pooled reversal is driven primarily by GPT and Grok; per-model label trajectories remain heterogeneous. As payoffs grow into standard integer ranges, the pooled language-model corpus becomes more cooperative, whereas the evolutionary reference becomes more dominated by defection (figure~\ref{fig:egtllm}b,c).

This direction is stable across the eighty evolutionary parameter settings reported in the appendix: the stationary AllD share is higher at the largest than at the smallest scale in every setting, and the rise is strictly monotonic in 61 of 80 settings. None reproduces the direction of the pooled empirical trajectory. \rev{From $\lam = 0.25$ upward the corpus resembles AllD in 0.26 to 0.31 of agent-games against 0.45 to 1.00 in the baseline, the direction a prosocial preference would produce, although the prompt effects discussed in Section~\ref{sec:discussion} could produce it too.} We therefore treat the comparison as a descriptive reference gap rather than as evidence for a particular mechanism. Evolutionary dynamics describe selection among rules, whereas the corpus describes fixed language-model agents.

\rev{Figure~\ref{fig:egtllm}c is therefore a constructed contrast:} the gap is a descriptive difference between selection among rules and fixed language-model agents. The language conditions also vary the stimulus, so this comparison does not identify a cultural or cognitive mechanism.

\begin{figure}[!htbp]
\centering
\includegraphics[width=\linewidth]{figures/f_egt_vs_llm.pdf}
\caption{\textbf{The evolutionary baseline and the corpus move in opposite directions.} (a) The baseline's stationary distribution over the four canonical rules across the ten payoff scales, drawn as a stacked band. (b) The corpus's strategy mix on the same axis, drawn identically. (c) The AllD share alone, both accounts on one axis, with the gap shaded. \rev{Panel (c) is a constructed contrast between two different objects, selection among rules in the baseline and fixed language-model agents in the corpus, so the gap is descriptive and not a test of one against the other.} Baseline at $Z = 100$, $\beta = 0.1$, $\epsilon = 0.05$, ten rounds; the corpus pooled over five models, five languages and four persona pairings.}
\label{fig:egtllm}
\end{figure}

\subsection{Agents condition more strongly on their own past actions}\label{sec:conditioning}

The evolutionary model rests on an explicit behavioural mechanism that can be evaluated directly in the data. In the classical account, cooperation survives invasion by defectors through \emph{reciprocity}: conditioning one's choice on what the partner just did, retaliating against defection and rewarding cooperation. Tit-for-Tat is the prototypical exemplar of this logic.

Because the previous round consists of an action pair (own action, partner action), the four conditional cooperation probabilities decompose cleanly into a base level and three orthogonal contrasts. \emph{Reciprocity} represents the partner main effect: the difference in cooperation after a partner cooperates versus after a partner defects. \emph{Persistence} represents the own main effect: the tendency to repeat one's own previous action. \emph{Matching} represents the interaction effect characteristic of Win-Stay-Lose-Shift. Textbook strategies occupy precise locations in this geometric space: Tit-for-Tat lies entirely on the reciprocity axis, Win-Stay-Lose-Shift on the matching axis, and unconditional rules at the origin (figure~\ref{fig:strategyspace}a).

Decomposing agent play reveals that own-action conditioning exceeds partner conditioning at the model level (figure~\ref{fig:strategyspace}). Across all five models, the persistence contrast is larger than the reciprocity contrast, although the margin is modest for Claude and substantially larger for the other models. Reciprocity remains positive in four models and is negative in Qwen. An unsupervised principal component analysis on raw transition vectors independently recovers this identical geometry without any prior knowledge of the contrasts (figure~\ref{fig:strategyspace}c).

The contrast estimates are precise at the dyad level. Whole-dyad bootstrap intervals for reciprocity are positive for Claude, GPT, Gemini~3.5, and Grok, and negative for Qwen; the corresponding persistence intervals are positive for all five models (electronic supplementary table~D.2). Thus the headline comparison is one of relative strength, not a claim that partner information is absent.

This finding narrows one candidate explanation for the divergence from the evolutionary baseline. In these data, agents condition more strongly on their own preceding moves than on their partner's behaviour. The contrast does not show that partner information is absent, nor does it identify the causal basis of cooperation.

A technical caution governs this decomposition. The naive difference between disagreement states, $p_{CD} - p_{DC}$, is frequently mistaken for a reciprocity metric in the literature. In reality, it equals persistence minus reciprocity identically, confounding the two main effects and falsely indicating anti-reciprocal behaviour where the true partner effect is positive.

\begin{figure}[!htbp]
\centering
\includegraphics[width=\linewidth]{figures/f_strategy_space.pdf}
\caption{\textbf{Own-action conditioning exceeds partner-action conditioning at the model level.} (a) Cells in the reciprocity by persistence plane, with the canonical rules at their exact coordinates; AllC, AllD and \rev{Win-Stay-Lose-Shift} all have both main effects zero and so coincide at the origin. (b) The three contrasts per model, against the value each canonical rule would give, with 95\% percentile intervals from whole-dyad bootstrap resampling. (c) The first two principal components of the raw five-vector, which is given no knowledge of the contrasts and recovers the same geometry. Panels (a) and (c) use model-by-scale-by-language cells; a predecessor observed fewer than thirty times yields no estimate and the cell is dropped, leaving 212 of 250.}
\label{fig:strategyspace}
\end{figure}


\section{Concluding remarks and discussion}\label{sec:discussion}

Positive scaling of payoffs is theoretically behaviour-preserving in the (repeated) prisoner's dilemma: \rev{it} does not alter the strategic structure of the game and therefore implies no change in predicted outcomes. Against this, our results show systematic departures from scale invariance across all five language models examined. These departures are not reducible to a homogeneous shift in baseline cooperation. Their magnitude varies across prompt languages, their direction and shape differ across models, and they are substantially attenuated by aggregation. The resulting aggregated trajectory is therefore not an adequate representation of any individual model's behavioural response. Relative to an evolutionary game baseline computed on the same payoff matrices, the aggregate corpus exhibits a decline in AllD-like behaviour as payoffs increase, contrary to the baseline prediction.
%This comparison quantifies the distance between the incentives encoded by the game and the behaviour produced by the agents, without by itself identifying the mechanism responsible for that distance.

The observed violation of scale invariance has important implications for the interpretation of cooperation benchmarks. A cooperation rate can no longer be regarded as a scale-independent property of a model, but must instead be understood as contingent on how the interaction is represented. This point resonates with the machine behaviour perspective, which argues that autonomous systems should be studied empirically, using the tools of the behavioural sciences, because their behaviour cannot be inferred from their construction alone~\cite{rahwan2019machine}. Yet behavioural experiments face a familiar difficulty: most manipulations do not imply a precise predicted effect, so any observed change must be interpreted relative to a comparison condition chosen by the experimenter. The invariance classes of decision theory provide an unusual exception. They specify manipulations for which the predicted behavioural effect is exactly zero. A violation of such an invariance is therefore especially informative: it does not merely show that behaviour changes under a different experimental condition, but that it changes under a transformation that should be behaviourally inert.

This point has consequences for several strands of machine-behaviour research touched by this study. Persona prompting is the clearest case in which the experimenter addresses the system in terms of its social role; our results show that whether such an instruction takes hold can depend on the numerical units in which the payoff matrix is presented. The implication is equally direct for models of mixed human--machine populations, which treat the cooperativeness of artificial agents as a parameter and ask where the population settles~\cite{santos2018social,han2013good,perc2017statistical,paiva2018engineering}. A cooperation rate partly determined by the experimenter's choice of units is not the scale-independent quantity such models require.


The methodological implication of our findings is significant. Cooperation rates elicited from language models are often treated as properties of the model under evaluation, conditional on a prompt and a game form. The present evidence indicates that they are instead joint properties of the model, the prompt language, the persona instruction, and the numerical scale at which the payoff matrix is expressed. The last of these is especially important because it is usually treated as theoretically irrelevant. In standard game-theoretic analysis, multiplying every payoff-matrix entry by a positive constant changes units but not incentives. For the models studied here, however, this formally irrelevant choice systematically changes observed behaviour. Thus, the payoff scale cannot be regarded as a neutral implementation detail in empirical studies of machine cooperation.

This conclusion is strengthened by the structure of the heterogeneity. The observed scale effects are not merely noisy fluctuations around an invariant response. They differ by model, interact with prompt language, moderate the behavioural effect of persona instructions, and alter the composition of inferred strategic patterns. These features are difficult to reconcile with sampling variability alone. They suggest that the models are sensitive to aspects of the numerical representation of the game that are not part of its strategic equivalence class. However, because the design manipulates payoff scale rather than the internal representation of number, notation, or stakes, the results should be interpreted as evidence of invariance violation rather than as evidence for any particular psychological or computational mechanism.

%The aggregation results are especially consequential. In the pooled data, a statistically significant main effect of payoff scale remains, but it explains substantially less variation than the model-by-scale interaction. Moreover, the pooled range understates the typical model-specific range by a factor of 1.62. The pooled mean therefore summarizes a behavioural trajectory that no individual model need follow. This is not simply a loss of precision that could be remedied by increasing the number of dyads. When model-specific scale responses differ in shape or direction, additional observations tighten confidence intervals around an aggregate that is itself substantively ambiguous. The appropriate unit of analysis is therefore the model-specific response surface, with prompt language and persona treated as factors that may themselves condition scale sensitivity.

%This point bears directly on benchmarking. A single cooperation rate measured at one payoff scale cannot be assumed to generalize across equivalent numerical representations of the same game. Two studies using the same model, prompt template, and game form could obtain markedly different cooperation rates if their payoff matrices differ only by a constant multiplier. Conversely, pooled comparisons across models may obscure behaviourally meaningful scale effects when models respond in opposing directions. Benchmark reports should therefore record the numerical scale of the payoff matrix and, where possible, include a scale-sensitivity analysis alongside any headline cooperation rate. Such analyses are particularly important when cooperation rates are subsequently used as inputs to population models, evolutionary simulations, or safety evaluations.

We also qualify the interpretation of persona effects. Persona instructions are often treated as relatively direct interventions on behavioural disposition, such as inducing more cooperative or more selfish play. The present results show that their effect may depend on the payoff scale at which the interaction is framed. In several LLMs, persona instructions reverse their behavioural effect below the unit threshold: under fractional stakes, models instructed to behave selfishly may cooperate more than those instructed to behave cooperatively. This pattern does not imply that the models possess stable preferences that are being inverted by small payoffs. Rather, it shows that the behavioural consequence of a verbal instruction is conditional on the numerical context in which that instruction is applied. Persona effects are therefore not separable from payoff representation.

The strategy-composition analysis extends the same conclusion beyond marginal cooperation rates. Payoff scale changes not only how often cooperation occurs, but also the kinds of action sequences that the models generate. These shifts indicate qualitative resemblances to canonical repeated-game strategies rather than exact implementations of them. The distinction is important. A ten-round trajectory cannot establish that a model is executing Tit-for-Tat, Win-Stay-Lose-Shift, Grim Trigger, or any other  behavioural strategy. What can be inferred is that some trajectories are more consistent with the behavioural signatures of such strategies than others, and that the prevalence of these signatures changes with payoff scale. Scale sensitivity therefore operates at the level of sequential structure as well as aggregate cooperation.

%The relationship between these findings and existing work on language-mediated cooperation is complementary. 

Prior studies have examined how cooperation changes under moral framing, pre-play communication, institutional mediation, or other contextual interventions \cite{buscemi2025strategic,lu2024llms,mao2025alympics,xie2026evolutionary}. Those manipulations can  alter the game as interpreted by the agent, because they change the semantic environment in which the decision is made. Payoff rescaling is different: under the standard formal model, it preserves the game. The fact that behaviour nevertheless changes indicates that language models may not respect the invariance classes assumed by game-theoretic representations. More broadly, if model behaviour is governed by a utility-like object that combines numerical payoffs with linguistic or representational features, then scale invariance will generally fail unless the numerical and non-numerical components transform together. Any theoretical account of machine play in strategic settings must therefore specify which transformations it regards as behaviourally equivalent.

\rev{These departures do not resemble human stake effects. Human effects are modest, tend to share a direction within a game and appear at stakes large relative to income; here a change of units with no consequence moves cooperation by up to 0.636, in opposite directions across models, and mostly at or below $\lam = 1$. Four explanations, not mutually exclusive, fit parts of the evidence: prompt text that does not scale with $\lam$, such as the narrative and the persona, consistent with the scale-dependent persona effects; the representation of numbers, since sub-unit payoffs are rare and tokenised differently, consistent with the change concentrating at or below $\lam = 1$; provider-specific post-training, consistent with Qwen and Grok inverting the persona at every scale; and a distorted human stake response learned from text. Appendix~F gives the evidence for each and a design that would test it. The present design cannot separate them, so the results identify a robust behavioural phenomenon but not its causal source.}


Several limitations follow from this design. First, the study examines five models, one repeated game, one family of prompt templates, and ten payoff scales. The results therefore support claims about the presence and structure of scale sensitivity in this corpus, not universal claims about all language models or all strategic environments. Second, the language comparison is bounded by the prompt templates used. Two templates describe payoffs as penalties while retaining a reward-maximization objective. Because this wording is fixed within each language across payoff scales, it does not undermine the within-language rescaling comparison. It may, however, interact with numerical magnitude and therefore limits any attribution of cross-language differences to language alone.
% Third, the statistical reference points used here quantify departures relative to the variability observed in the corpus. The permutation tests and dyad-level bootstrap establish that the observed ranges are unlikely under the relevant exchangeability assumptions, but they do not provide a full hierarchical model of replicate, dyad, prompt, and game-level dependence. Clustering at the dyad level addresses important dependence among actions within repeated interactions, yet a more complete design would include independent replications and mixed-effects models with random intercepts or slopes for replicate and game. The independent repeat of Gemini 3.5 at four scales, reported in the appendix, provides useful evidence on repeatability for that arm, but it should not be interpreted as a general test-retest estimate for the full corpus.
Third, the strategy analysis is necessarily approximate. Strategies are latent decision rules, whereas the data consist of finite action sequences. Ten rounds provide limited information for distinguishing among memory-one rules, especially when execution noise is present. The classifier was validated up to 20\% execution noise, corresponding to two deviations in a ten-round sequence, and some models operate near that boundary on average. The resulting labels should therefore be read as descriptive summaries of trajectory resemblance, not as evidence that the models implement the corresponding strategies. Longer horizons, richer strategy classes, and mixture-based estimators would be needed to make stronger claims about strategic implementation \cite{inferenceStrategies,hilbe2013memory,han2013intention}. 

\rev{Fourth, there is no condition without numbers. A verbal baseline describing each outcome only by its rank would keep the ordinal structure while removing every magnitude, and would show whether numerical payoffs themselves change behaviour (appendix~F.3).}
%Last but not least, the evolutionary baseline should be interpreted as a theoretical reference distribution rather than as a fitted model of language-model behaviour. It maps stochastic behaviour onto a symmetric per-action error process, which is analytically useful but necessarily coarse. Prompt sensitivity, model heterogeneity, language effects, end-game reasoning, and scale-dependent instruction following are not well captured by symmetric action flips. The baseline is therefore valuable precisely because it clarifies the direction and magnitude of the discrepancy between formal incentive predictions and observed play, not because it offers a complete generative account of the data. 

Overall, our study shows that payoff scale, although formally irrelevant under standard game-theoretic invariance, systematically affects machine cooperation. The effect is model-specific, language-dependent, and consequential for both cooperation rates and inferred strategic composition. These findings call for caution in interpreting single-scale cooperation benchmarks and suggest that scale sensitivity should be reported as a standard robustness measure in studies of language models in games. More generally, they indicate that empirical work on machine strategic behaviour cannot assume that formal equivalences in the game are automatically respected by the agents that are prompted to play it.


%%%%%%%%%% End matter %%%%%%%%%%

\section*{Ethics statement}
This work did not require ethical approval from a human subject or animal welfare committee.

\section*{Data accessibility}
The game transcripts, prompt templates, trained strategy classifier and analysis code supporting this study are \rev{deposited at Zenodo, \url{https://doi.org/10.5281/zenodo.23627009}~\cite{han2026data}. The deposit holds the per-game tables of all 12{,}000 dyads, the prompts in all five languages, the raw collection records, the game code, the classifier and the analysis pipeline that regenerates every table and figure.}

\section*{Declaration of AI use}
Large language models are the object of study in this article. Generative AI was used only as an assistant to identify and fix bugs in code and to edit English grammar.

\section*{Author contributions}
T.A.H.: conceptualisation, methodology, supervision, project administration, writing (original draft), writing (review and editing). T.-K.H.: conceptualisation, software, formal analysis, visualisation, writing (original draft). D.-S.D.-M.: software, data curation, formal analysis, writing (original draft). T.-B.C., P.-H.L., H.-D.N., P.-Q.N.-L., M.-L.N.-V., H.-P.P., P.-H.P., T.-K.T., C.-N.T., H.T. and G.-T.T.-L.: investigation, data curation, validation, writing (review and editing). A.B.: methodology, resources, writing (review and editing). L.H.T.: conceptualisation,  supervision,  writing (review and editing). T.A.H., T.-K.H. and D.-S.D.-M. contributed equally to this work. All authors gave final approval for publication and agree to be held accountable for the work performed therein.

\section*{Competing interests}
We declare we have no competing interests.

\section*{Funding}
T.A.H. is supported by EPSRC (grant EP/Y00857X/1) and also acknowledges travel support from HCMUT\text{-}VNUHCM.

%%%%%%%%%% Bibliography %%%%%%%%%%
\bibliographystyle{unsrtnat}
\bibliography{mybib}

\end{document}
